\subsection{Structure of indecomposable injective modules}

Lemma 1.--- Let A be a ring, with a an ideal of A, and I an A-module. For every integer \(n \geqslant 0\) let us denote by \(I_{n}\) the submodule of I consisting of the elements annihilated by \(a^{n}\).

\begin{enumerate}
\def\labelenumi{\alph{enumi})}
\item
  Suppose the A-module I is injective. Then the \(\mathrm{A} / \mathfrak{a}^n\)-module \(\mathrm{I}_n\) is injective for every \(n \geqslant 0\).
\item
  Suppose that the ring A is Noetherian, and that for every \(n \geqslant 0\) the \(\mathrm{A} / \mathfrak{a}^n\) -module \(\mathrm{I}_n\) be injective. Then the union of the \(\mathrm{I}_n\) is an injective A-module.
\item
  The \(\mathbf{A} / \mathfrak{a}^n\) -module \(\mathrm{I}_n\) is isomorphic to \(\operatorname{Hom}_{\mathbf{A}}(\mathbf{A} / \mathfrak{a}^n, \mathbf{I})\), which is injective (A, X, p. 18, prop. 11).
\item
  Let J be the union of the \(I_{n}\), b an ideal of A, and \(f : b \to J\) an A-homomorphism. It remains (A, X, p. 16, prop. 10) of proving that there exists an element x of J such that one has \(f(b) = bx\) for every \(b \in b\). Since b is of finite type, there exists an integer n such that one has \(f(\mathfrak{b}) \subset \mathrm{I}_{n}\) , that is \(f(\mathfrak{a}^{n}\mathfrak{b}) = 0\) . According to cor. 2 of prop. 1 of III, § 3, no. 1, there exists an integer \(m \geqslant n\) such that \(a^{m} \cap b \subset a^{n}b\) , hence \(f(\mathfrak{a}^{m} \cap \mathfrak{b}) = 0\) . Then f induces an A/\(a^{m}\)-linear homomorphism from \(\mathfrak{b}/(\mathfrak{a}^{m} \cap \mathfrak{b})\) into \(I_{m}\); since the A/ \(a^{m}\)-module \(I_{m}\) is injective, there exists an element x of \(I_{m}\) such that one has \(f(b) = bx\) for every \(b \in b\) , whence b).
\end{enumerate}

Let \(\mathfrak{a}\) be an ideal of A; in what follows we shall agree to set \(\mathfrak{a}^n = \mathrm{A}\) for every integer \(n \leqslant 0\). Let E be an A-module. For every \(n \in \mathbf{Z}\) let us denote by \(\mathrm{E}_n\) the submodule of E consisting of the elements annihilated by \(\mathfrak{a}^n\); let \(\mathrm{gr}^{\mathfrak{a}}(\mathrm{E})\) be the graded A-module of type \(\mathbf{Z}\) such that \(\mathrm{gr}^{\mathfrak{a}}(\mathrm{E})_m = \mathrm{E}_{-m+1}/\mathrm{E}_{-m}\) for every integer \(m\). The module \(\mathrm{gr}^{\mathfrak{a}}(\mathrm{E})_m\) is zero for \(m \geqslant 1\) , and \(\mathrm{gr}^{\mathfrak{a}}(\mathrm{E})_0\) is identified with \(\mathrm{E}_1\) . Let \(\mathrm{gr}(\mathrm{A})\) be the associated graded ring of A for the \(\mathfrak{a}\) -adic filtration: we have \(\mathrm{gr}(\mathrm{A})_n = \mathfrak{a}^n/\mathfrak{a}^{n+1}\) for every \(n \in \mathbf{Z}\). Let \(n\) and \(m\) be integers. By passing to quotients of the bilinear map \((a,x) \mapsto ax\) of \(\mathfrak{a}^n \times \mathrm{E}_{-m+1}\) in \(\mathrm{E}_{-m-n+1}\) we deduce an A/ \(\mathfrak{a}\) -bilinear map

\[
\alpha_ {n, m}: \mathrm{gr} (\mathrm{A}) _ {n} \times \mathrm{gr} ^ {\mathfrak {a}} (\mathrm{E}) _ {m} \longrightarrow \mathrm{gr} ^ {\mathfrak {a}} (\mathrm{E}) _ {n + m},\tag{1}
\]

which defines on gr \(^{a}\) (E) a graded gr(A)-module structure. For every \(n \in \mathbf{Z}\) one deduces from the A/a-bilinear map \(\alpha_{n,-n} : \text{gr}(A)_n \times \text{gr}^{a}(E)_{-n} \longrightarrow E_1\) an A/a-linear map \(\beta_{E,n} : \text{gr}^{a}(E)_{-n} \longrightarrow \text{Hom}_{A/\alpha}(\text{gr}(A)_n, E_1)\); the maps \(\beta_{E,n}\) are the components of a homomorphism of graded A/a-modules, called canonical.

\[
\beta_ {\mathrm{E}}: \operatorname{gr} ^ {a} (\mathrm{E}) \longrightarrow \operatorname{Homgr} _ {\mathrm{A} / \mathfrak{a}} (\operatorname{gr} (\mathrm{A}), \mathrm{E} _ {1}).\tag{2}
\]

For \(a \in \text{gr(A)}\), \(x \in \text{gr}^{\mathfrak{a}}(\text{E})\) , \(\beta_{\text{E}}(x)(a)\) is by definition the component in \(\text{gr}^{\mathfrak{a}}(\text{E})_0 = \text{E}_1\) of the element \(ax\) of \(\text{gr}^{\mathfrak{a}}(\text{E})\). It follows that \(\beta_{\text{E}}\) is \(\text{gr(A)}\) -linear when one equips \(\text{Homgr}_{\text{A}/\mathfrak{a}}(\text{gr(A)}, \text{E}_1)\) with the structure of \(\text{gr(A)}\) -module defined by the formula \((bf)(a) = f(ab)\) for \(a, b\) in \(\text{gr(A)}\) and \(f\) in \(\text{Homgr}_{\text{A}/\mathfrak{a}}(\text{gr(A)}, \text{E}_1)\) .

PROPOSITION 2.--- Let A be a Noetherian ring, with a an ideal of A, E an A-module, and M a sub-A-module of E annihilated by a. The following conditions are equivalent:

\begin{enumerate}
\def\labelenumi{(\roman{enumi})}
\item
  E is an injective envelope of M;
\item
  the \(\mathrm{A}/\mathfrak{a}\)-module \(\mathrm{E}_1\) is an injective envelope of the \(\mathrm{A}/\mathfrak{a}\) -module \(\mathrm{M}\); the module \(\mathrm{E}\) is the union of \(\mathrm{E}_n\), and the canonical map \(\beta_{\mathrm{E}}\) is bijective.
\end{enumerate}

Suppose condition (i) is satisfied. The A/a-module \(E_{1}\) is injective (Lemma 1, a)), and contains M; since every sub- \(A/a\) -module of \(E_{1}\) is a sub-A-module of E, \(E_{1}\) is an injective envelope of the A/a-module M. By Lemma 1, the union of the \(E_{n}\) is an injective sub-A-module of E containing M, hence equal to E. Since E is injective, we have for all \(n \geqslant 0\) an exact sequence

\[
0 \to \operatorname{Hom} _ {\mathrm{A}} (\mathrm{A} / \mathfrak {a} ^ {n}, \mathrm{E}) \longrightarrow \operatorname{Hom} _ {\mathrm{A}} (\mathrm{A} / \mathfrak {a} ^ {n + 1}, \mathrm{E}) \longrightarrow \operatorname{Hom} _ {\mathrm{A}} (\mathfrak {a} ^ {n} / \mathfrak {a} ^ {n + 1}, \mathrm{E}) \to 0;
\]

as \(\mathrm{Hom}_{\mathrm{A}}(\mathrm{A}/\mathfrak{a}^{m},\mathrm{E})\) is identified with \(\mathrm{E}_{m}\) for every \(m\) and that the canonical injection of \(\mathrm{Hom}_{\mathrm{A}}(\mathfrak{a}^{n}/\mathfrak{a}^{n+1},\mathrm{E}_{1})\) in \(\mathrm{Hom}_{\mathrm{A}}(\mathfrak{a}^{n}/\mathfrak{a}^{n+1},\mathrm{E})\) is bijective, we deduce that the canonical homomorphism \(\beta_{\mathrm{E}}\) is bijective, whence (ii).

Suppose (ii) is satisfied. Let \(e: \mathrm{M} \to \mathrm{I}\) be an injective envelope of M. Since I is injective, there exists an A-linear map \(\varphi: \mathrm{E} \to \mathrm{I}\) extending \(e\). But \(\varphi\) maps \(\mathrm{E}_n\) in \(\mathrm{I}_n\) for every \(n\), therefore induces homomorphisms \(\mathrm{gr}^{\mathfrak{a}}(\varphi):\mathrm{gr}^{\mathfrak{a}}(\mathrm{E})\to \mathrm{gr}^{\mathfrak{a}}(\mathrm{I})\) and \(\varphi_{1}:\mathrm{E}_{1}\rightarrow \mathrm{I}_{1}\) thereby making the diagram commutative

\[
\begin{array}{c c c} \operatorname{gr} ^ {\mathfrak {a}} (\mathrm{E}) & \xrightarrow {\beta_ {\mathrm{E}}} & \operatorname{Homgr} _ {\mathrm{A} / \mathfrak {a}} (\operatorname{gr} (\mathrm{A}), \mathrm{E} _ {1}) \\ \operatorname{gr} ^ {\mathfrak {a}} (\varphi) \Bigg \downarrow & & \Bigg \downarrow \operatorname{Homgr} (1, \varphi_ {1}) \\ \operatorname{gr} ^ {\mathfrak {a}} (\mathrm{I}) & \xrightarrow {\beta_ {\mathrm{I}}} & \operatorname{Homgr} _ {\mathrm{A} / \mathfrak {a}} (\operatorname{gr} (\mathrm{A}), \mathrm{I} _ {1}). \end{array}
\]

Since \(E_{1}\) and \(I_{1}\) are injective envelopes of the A/a-module M, the homomorphism \(\varphi_{1}\) is bijective; since \(\beta_{E}\) and \(\beta_{I}\) are bijective, it follows that \(\mathrm{gr}^{\mathfrak{a}}(\varphi)\) is bijective. This implies, by induction on n, that \(\varphi\) induces a bijection from \(E_{n}\) onto \(I_{n}\) for every \(n \geqslant 1\); therefore \(\varphi\) is bijective, which implies (i).

Lemma 2.--- Let A be a ring, let a be a finitely generated ideal of A, and let M be an A-module in which every element is annihilated by a power of a. Let A-hat denote the separated completion of A for the a-adic topology.

\begin{enumerate}
\def\labelenumi{\alph{enumi})}
\item
  There exists on M a unique structure of \(\widehat{\mathrm{A}}\) -module extending the given structure of A-module.
\item
  The sub-A-hat-modules of M are its sub-A-modules, and one has Hom \(_{\text{A}}\) (M, P) = Hom \(_{\widehat{\text{A}}}\) (M, P) for every \(\widehat{\text{A}}\) -module P.
\item
  Identify \(\widehat{\mathbf{A}}\) with the projective limit of the rings \(\mathrm{A} / \mathfrak{a}^n\), and endow M with the discrete topology. Let \(a = (a_n)\) be an element of \(\widehat{\mathbf{A}}\) , and let \(x\) an element of M. Since \(x\) is annihilated by a power of \(\mathfrak{a}\) , the sequence \((a_n x)\) is stationary; denote its limit by \(ax\). The map \((a, x) \mapsto ax\) defines on M a structure of \(\widehat{\mathbf{A}}\) -module which extends the given structure of A-module.
\end{enumerate}

Conversely, suppose such a structure on M is given; let \(a = (a_{n})\) be an element of \(\widehat{\mathrm{A}}\), \(x\) be an element of M and \(m\) an integer such that \(\mathfrak{a}^{m}x = 0\). For every integer \(n\) , \(a - a_{n}\) belongs to \(\widehat{\mathfrak{a}^{n}}\) , which is equal to \(\mathfrak{a}^{n}\widehat{\mathrm{A}}\) (III, § 2, n° 12, cor. 2 of prop. 16); one therefore has \(ax = a_{n}x\) for \(n \geqslant m\) , whence the uniqueness assertion.

\begin{enumerate}
\def\labelenumi{\alph{enumi})}
\setcounter{enumi}{1}
\tightlist
\item
  It follows from the preceding that one has \(Ax = \widehat{A}x\) for every \(x \in M\) ; the sub- \(\widehat{A}\)-modules of \(M\) are therefore its sub- \(A\) -modules. Finally, let \(u\) be an \(A\)-linear homomorphism from \(M\) into a \(\widehat{A}\) -module \(P\). Let \(a = (a_n)\) be an element of \(\widehat{A}\) , \(x\) be an element of \(M\) and \(m\) an integer such that \(\mathfrak{a}^m x = 0\) ; one has \(\mathfrak{a}^m u(x) = 0\) . Since \(a - a_m\) belongs to \(\mathfrak{a}^m\widehat{A}\) , one has
\end{enumerate}

\[
u (a x) = u \left(a _ {m} x\right) = a _ {m} u (x) = a u (x),
\]

so that u is A-hat-linear.

PROPOSITION 3. Let A be a Noetherian ring, p a prime ideal of A and e : A/p → I an injective envelope of the A-module A/p. For every integer n ≥ 0, denote by I \(_{n}\) the submodule of I consisting of the elements annihilated by p \(^{n}\).

\begin{enumerate}
\def\labelenumi{\alph{enumi})}
\item
  The A-module I is the union of the \(I_{n}\). The injection \(\Lambda/p \to I_{1}\) extends to an isomorphism from \(\kappa(\mathfrak{p})\) onto \(I_{1}\); identify \(\kappa(\mathfrak{p})\) with \(I_{1}\) by means of this isomorphism. For each integer \(n \geqslant 0\) , the structure of \(\Lambda/p\)-module on \(I_{n+1}/I_{n}\) comes by restriction of scalars from a unique structure of \(\kappa(\mathfrak{p})\) -vector space; the canonical homomorphism \(\beta_{I,-n}: I_{n+1}/I_{n} \longrightarrow \operatorname{Hom}_{\mathrm{A}/\mathfrak{p}}(\mathfrak{p}^{n}/\mathfrak{p}^{n+1}, \kappa(\mathfrak{p}))\) is an isomorphism of \(\kappa(\mathfrak{p})\) -vector spaces of finite dimension.
\item
  There exists a unique structure of \(\mathrm{A_p}\) -module on I inducing its structure of A-module. The canonical homomorphism \(\widehat{\mathrm{A_p}} \longrightarrow \mathrm{End}_{\mathrm{A}}(\mathrm{I})\) is bijective.
\end{enumerate}

By A, X, p. 20, example 1, the A/p-module \(\kappa(\mathfrak{p})\) is an injective envelope of A/p. It therefore follows from prop. 2 that I \(_1\) is identified with \(\kappa(\mathfrak{p})\), that I is the union of the I \(_m\) , and for every integer \(n \geqslant 0\) , \(\beta_{1,-n}\) is an isomorphism of A/p-modules. For every nonzero element \(a\) of A/p, the homothety with ratio \(a\) is invertible in Hom \(_{\text{A/p}}(\mathfrak{p}^n/\mathfrak{p}^{n+1}, \kappa(\mathfrak{p}))\) , hence also in I \(_{n+1}\) /I \(_n\) , which completes the proof of a).

Let \(s \in A - p\). Since the homothety \(s_{A/p}\) is injective, the trace of Ker \(s_I\) on \(A/p\) is zero, which implies that the homothety \(s_I\) is injective. Then \(sI\) is a direct summand of I (A, X, p. 19, cor. 4), hence equal to I since I is indecomposable ( \(n^o\) 1, prop. 1), so that the homothety \(s_I\) is bijective. There therefore exists a unique structure of \(A_p\) -module on I inducing its structure of A-module; it extends uniquely to a structure of \(\widehat{A_p}\) -module (lemma 2).

For every integer \(n\) , we deduce from the canonical ring homomorphism \(\Lambda_{\mathfrak{p}} \longrightarrow \mathrm{End}_{\mathrm{A}}(\mathrm{I})\) an A-linear map \(\alpha_{n}: \mathrm{A}_{\mathfrak{p}}/\mathfrak{p}^{n}\mathrm{A}_{\mathfrak{p}} \longrightarrow \mathrm{Hom}_{\mathrm{A}}(\mathrm{I}_{n},\mathrm{I})\). Consider the commutative diagram with exact rows

\[
\begin{array}{c c c c c c c c c} 0 & \longrightarrow & \mathfrak {p} ^ {n} A _ {\mathfrak {p}} / \mathfrak {p} ^ {n + 1} A _ {\mathfrak {p}} & \longrightarrow & A _ {\mathfrak {p}} / \mathfrak {p} ^ {n + 1} A _ {\mathfrak {p}} & \longrightarrow & A _ {\mathfrak {p}} / \mathfrak {p} ^ {n} A _ {\mathfrak {p}} & \longrightarrow & 0 \\ & & \alpha_ {n + 1} ^ {\prime} \Bigg \downarrow & & \alpha_ {n + 1} \Bigg \downarrow & & \alpha_ {n} \Bigg \downarrow \\ 0 & \longrightarrow & \operatorname{Hom} _ {A} (I _ {n + 1} / I _ {n}, I _ {1}) & \longrightarrow & \operatorname{Hom} _ {A} (I _ {n + 1}, I) & \longrightarrow & \operatorname{Hom} _ {A} (I _ {n}, I) & \longrightarrow & 0 \end{array}
\]

where \(\alpha_{n+1}^{\prime}\) is the homomorphism induced by \(\alpha_{n+1}\). Consider the canonical \(\kappa(\mathfrak{p})\)-bilinear map

\[
\alpha_ {n, - n}: \mathfrak {p} ^ {n} \mathrm{A} _ {\mathfrak {p}} / \mathfrak {p} ^ {n + 1} \mathrm{A} _ {\mathfrak {p}} \times \mathrm{I} _ {n + 1} / \mathrm{I} _ {n} \longrightarrow \mathrm{I} _ {1}
\]

(formula (1)). The linear map \(I_{n+1}/I_{n} \longrightarrow \operatorname{Hom}_{\kappa(\mathfrak{p})}(\mathfrak{p}^{n}A_{\mathfrak{p}}/\mathfrak{p}^{n+1}A_{\mathfrak{p}}, I_{1})\) associated with it on the left is identified with \(\beta_{I,-n}\), and the one associated with it on the right is \(\alpha_{n+1}'\) . Since \(\beta_{I,-n}\) is bijective by a), and the same holds for \(\alpha_{n+1}'\); it follows from the above diagram, by induction on n, that \(\alpha_{n}\) is an isomorphism for every n. Since I is the union of the \(I_{n}\) , the canonical map \(\operatorname{End}_{\mathrm{A}}(\mathrm{I}) \to \varprojlim \operatorname{Hom}_{\mathrm{A}}(\mathrm{I}_{n}, \mathrm{I})\) is bijective; the ring homomorphism \(\widehat{A_{p}} \to \operatorname{End}_{\mathrm{A}}(\mathrm{I})\) which is identified with the projective limit of the maps \(\alpha_{n}\) is therefore bijective.

Remark.--- It follows from the preceding proof that the annihilator of \(I_{n}\) in \(\widehat{A_{p}}\) (resp. in \(A_{p}\) ) is \(p^{n}\widehat{A_{p}}\) (resp. \(p^{n}A_{p}\) ). Consequently, the annihilator of the A-module \(I_{n}\) is the inverse image in A of the ideal \(p^{n}A_{p}\) which is sometimes denoted \(\mathfrak{p}^{(n)}\) which is called the n-th symbolic power of the prime ideal p.

COROLLARY.--- Let J be an injective A-module such that \(\mathrm{Ass}_{\mathrm{A}}(\mathrm{J}) = \{\mathfrak{p}\}\).

\begin{enumerate}
\def\labelenumi{\alph{enumi})}
\item
  \(L^{\prime}\) canonical map \(\mathrm{J} \to \mathrm{A}_{\mathfrak{p}} \otimes_{\mathrm{A}} \mathrm{J}\) is bijective.
\item
  Let E denote the A/p-module Hom \(_A\) (A/p, J). On E there exists a unique structure of \(\kappa(p)\)-vector space extending its structure as an A/p-module; the A-module J is isomorphic to I \(^{([E:\kappa(p)])}\).
\end{enumerate}

Indeed, J is isomorphic to an A-module \(I^{(\mathfrak{c})}\) , where c is a suitable cardinal (no. 1, th. 1). The corollary follows from the proposition when J = I, and the general case is deduced from it immediately.

